% GENERATED FILE. Edit canonical lessons, then run npm run generate:books.
% canonical-source-hash: fnv1a64:18f5290f82438989
% canonical-lessons: LA-C486-pruina, LA-C486-frigus, LA-C486-gelo, LA-C486-rigeo, LA-C486-algeo, LA-R486-review-479-486

\chapter{Hoarfrost, The cold, To freeze, To be stiff with cold, To feel cold}
\label{ch:la-hoarfrost-the-cold-to-freeze-to-be-stiff-with-cold-to-feel-cold}
\hlchaptermodality{486}

\begin{chapteropening}
\textbf{By the end of this chapter:} \emph{I can talk about cold weather: hoarfrost, the cold, water freezing, being stiff with cold, and feeling cold.}

Complete the last lesson of chapter 486: I can talk about cold weather: hoarfrost, the cold, water freezing, being stiff with cold, and feeling cold.
\end{chapteropening}

\section[pruīna, pruīnae]{pruīna, pruīnae — hoarfrost, frost}

\label{lesson:LA-C486-pruina}



\begin{quote}
Before the new one: say the Latin for stormy, then the Latin for turbulent.
\end{quote}

\subsection*{You'll want to know: pruīna, pruīnae}
\textbf{pruīna, pruīnae} — \textquotedblleft{}hoarfrost, frost\textquotedblright{}.

\begin{quote}
Mane in prātīs pruīna alba est.
\end{quote}

\begin{quote}
\emph{In the morning there is white frost on the meadows.}
\end{quote}

First declension, feminine: \emph{pruīna, pruīnae f.} In poetry, also winter.

\begin{cousinweb}
\emph{pruīna} is inherited from PIE \emph{prews-} 'to frost, to freeze' (EtymDB). English \emph{freeze} and \emph{frost} are Germanic cousins from the same root.
\end{cousinweb}

\begin{grammarlens}[title={What you've built}]
One more word for telling, explaining and describing at B1.
\end{grammarlens}

\subsection*{Your turn}
\begin{itemize}
\raggedright
  \item \emph{Say it:} \emph{pruīna, pruīnae}
  \item \emph{Say it:} \emph{pruīna, pruīnae}, once more
  \item \emph{Say it:} a sentence of your own with \emph{pruīna, pruīnae}
  \item \emph{Recall:} say the Latin for stormy, then the Latin for turbulent, then say \emph{pruīna, pruīnae} again
\end{itemize}

\begin{tcolorbox}[breakable,colback=teal!4,colframe=teal!35!black,title={Before you move on}]
What does pruīna, pruīnae mean? (\textquotedblleft{}Hoarfrost, frost\textquotedblright{}.) Say it once more.
\end{tcolorbox}
\section[frīgus, frīgoris]{frīgus, frīgoris — the cold, cold weather}

\label{lesson:LA-C486-frigus}



\begin{quote}
Before the new one: say the Latin for turbulent, then the Latin for hoarfrost.
\end{quote}

\subsection*{You'll want to know: frīgus, frīgoris}
\textbf{frīgus, frīgoris} — \textquotedblleft{}the cold, cold weather\textquotedblright{}.

\begin{quote}
Hieme frīgus magnum est, et pauperēs ignem nōn habent.
\end{quote}

\begin{quote}
\emph{In winter the cold is great, and the poor have no fire.}
\end{quote}

Third declension, neuter: \emph{frīgus, frīgoris n.} The adjective \emph{frīgidus} 'cold' gives English \emph{frigid}.

\begin{cousinweb}
\emph{frīgus} is inherited from PIE \emph{sriHgos} 'cold, coldness' (EtymDB), with Greek \emph{rhigos} 'cold' as a cousin. English \emph{frigid} comes from \emph{frīgidus}, built on it (Lewis \& Short).
\end{cousinweb}

\begin{grammarlens}[title={What you've built}]
One more word for telling, explaining and describing at B1.
\end{grammarlens}

\subsection*{Your turn}
\begin{itemize}
\raggedright
  \item \emph{Say it:} \emph{frīgus, frīgoris}
  \item \emph{Say it:} \emph{frīgus, frīgoris}, once more
  \item \emph{Say it:} a sentence of your own with \emph{frīgus, frīgoris}
  \item \emph{Recall:} say the Latin for turbulent, then the Latin for hoarfrost, then say \emph{frīgus, frīgoris} again
\end{itemize}

\begin{tcolorbox}[breakable,colback=teal!4,colframe=teal!35!black,title={Before you move on}]
What does frīgus, frīgoris mean? (\textquotedblleft{}The cold, cold weather\textquotedblright{}.) Say it once more.
\end{tcolorbox}
\section[gelō, gelāre]{gelō, gelāre — to freeze; (passive) to be frozen}

\label{lesson:LA-C486-gelo}



\begin{quote}
Before the new one: say the Latin for hoarfrost, then the Latin for the cold.
\end{quote}

\subsection*{You'll want to know: gelō, gelāre}
\textbf{gelō, gelāre} — \textquotedblleft{}to freeze; (passive) to be frozen\textquotedblright{}.

\begin{quote}
Nocte aqua in puteō gelātur.
\end{quote}

\begin{quote}
\emph{At night the water in the well freezes.}
\end{quote}

\emph{gelō, gelāre, gelāvī, gelātum}. English \emph{congeal}, \emph{jelly} and \emph{gelatin} come from this family.

\begin{cousinweb}
\emph{gelō} is built on \emph{gelū} 'frost, chill', from PIE \emph{gel-} 'cold' (EtymDB; Lewis \& Short). English \emph{congeal} comes from \emph{congelāre}, and \emph{jelly} from \emph{gelāta} 'frozen' through French.
\end{cousinweb}

\begin{grammarlens}[title={What you've built}]
One more word for telling, explaining and describing at B1.
\end{grammarlens}

\subsection*{Your turn}
\begin{itemize}
\raggedright
  \item \emph{Say it:} \emph{gelō, gelāre}
  \item \emph{Say it:} \emph{gelō, gelāre}, once more
  \item \emph{Say it:} a sentence of your own with \emph{gelō, gelāre}
  \item \emph{Recall:} say the Latin for hoarfrost, then the Latin for the cold, then say \emph{gelō, gelāre} again
\end{itemize}

\begin{tcolorbox}[breakable,colback=teal!4,colframe=teal!35!black,title={Before you move on}]
What does gelō, gelāre mean? (\textquotedblleft{}To freeze; (passive) to be frozen\textquotedblright{}.) Say it once more.
\end{tcolorbox}
\section[rigeō, rigēre]{rigeō, rigēre — to be stiff with cold, to be numb; to be stiff}

\label{lesson:LA-C486-rigeo}



\begin{quote}
Before the new one: say the Latin for the cold, then the Latin for to freeze.
\end{quote}

\subsection*{You'll want to know: rigeō, rigēre}
\textbf{rigeō, rigēre} — \textquotedblleft{}to be stiff with cold, to be numb; to be stiff\textquotedblright{}.

\begin{quote}
Manūs viātōrum frīgore rigent.
\end{quote}

\begin{quote}
\emph{The travellers' hands are stiff with cold.}
\end{quote}

\emph{rigeō, rigēre, riguī}, second conjugation, no supine. English \emph{rigid} comes from \emph{rigidus}.

\begin{cousinweb}
The origin of \emph{rigeō} is not settled: Lewis \& Short suggest a link with Greek \emph{rhigeō} 'to shiver' and \emph{frīgeō}, and EtymDB gives no parent. English \emph{rigid} comes from \emph{rigidus}, built on it.
\end{cousinweb}

\begin{grammarlens}[title={What you've built}]
One more word for telling, explaining and describing at B1.
\end{grammarlens}

\subsection*{Your turn}
\begin{itemize}
\raggedright
  \item \emph{Say it:} \emph{rigeō, rigēre}
  \item \emph{Say it:} \emph{rigeō, rigēre}, once more
  \item \emph{Say it:} a sentence of your own with \emph{rigeō, rigēre}
  \item \emph{Recall:} say the Latin for the cold, then the Latin for to freeze, then say \emph{rigeō, rigēre} again
\end{itemize}

\begin{tcolorbox}[breakable,colback=teal!4,colframe=teal!35!black,title={Before you move on}]
What does rigeō, rigēre mean? (\textquotedblleft{}To be stiff with cold, to be numb; to be stiff\textquotedblright{}.) Say it once more.
\end{tcolorbox}
\section[algeō, algēre]{algeō, algēre — to feel cold, to be cold}

\label{lesson:LA-C486-algeo}



\begin{quote}
Before the new one: say the Latin for to freeze, then the Latin for to be stiff with cold.
\end{quote}

\subsection*{You'll want to know: algeō, algēre}
\textbf{algeō, algēre} — \textquotedblleft{}to feel cold, to be cold\textquotedblright{}.

\begin{quote}
Servus sine palliō alget.
\end{quote}

\begin{quote}
\emph{The slave without a cloak feels cold.}
\end{quote}

\emph{algeō, algēre, alsī}, second conjugation. Of a person feeling cold; \emph{frīgus} is the cold itself.

\begin{cousinweb}
The origin of \emph{algeō} is not settled: Lewis \& Short report an old link with Greek \emph{algeō} 'to feel pain', and EtymDB gives no parent.
\end{cousinweb}

\begin{grammarlens}[title={What you've built}]
One more word for telling, explaining and describing at B1.
\end{grammarlens}

\subsection*{Your turn}
\begin{itemize}
\raggedright
  \item \emph{Say it:} \emph{algeō, algēre}
  \item \emph{Say it:} \emph{algeō, algēre}, once more
  \item \emph{Say it:} a sentence of your own with \emph{algeō, algēre}
  \item \emph{Recall:} say the Latin for to freeze, then the Latin for to be stiff with cold, then say \emph{algeō, algēre} again
\end{itemize}

\begin{tcolorbox}[breakable,colback=teal!4,colframe=teal!35!black,title={Before you move on}]
What does algeō, algēre mean? (\textquotedblleft{}To feel cold, to be cold\textquotedblright{}.) Say it once more.
\end{tcolorbox}
\section[(dialogue)]{Words from chapters 479–486 — a review}

\label{lesson:LA-R486-review-479-486}



\begin{quote}
Say the Latin for an orchard and an olive tree.
\end{quote}

\subsection*{Your turn}
Say these aloud:

\begin{itemize}
\raggedright
  \item an orchard, an olive tree, a press, to graft, the grape harvest
  \item a cliff, a jagged rock, sheer, steep, to climb
  \item an enclosed valley, a plain, rocky, rough, to slide down
  \item a river, a torrent, a confluence, to flow, to swim across
  \item a pool, a whirlpool, a chasm, to trickle, to flood
  \item a beech tree, a fir tree, an elm tree, an ash tree, shady
  \item a downpour, to be soaking wet, rainy, stormy, turbulent
  \item hoarfrost, the cold, to freeze, to be stiff with cold, to feel cold
\end{itemize}

\begin{tcolorbox}[breakable,colback=teal!4,colframe=teal!35!black,title={Before you move on}]
An orchard? (\textbf{pōmārium, pōmāriī}.) An olive tree? (\textbf{olea, oleae}.) A press? (\textbf{prēlum, prēlī}.) To graft? (\textbf{īnserō, īnserere}.) The grape harvest? (\textbf{vindēmia, vindēmiae}.) A cliff? (\textbf{rūpēs, rūpis}.) A jagged rock? (\textbf{cautēs, cautis}.) Sheer? (\textbf{praeruptus, praerupta, praeruptum}.) Steep? (\textbf{arduus, ardua, arduum}.) To climb? (\textbf{scandō, scandere}.) An enclosed valley? (\textbf{convallis, convallis}.) A plain? (\textbf{plānitiēs, plānitiēī}.) Rocky? (\textbf{saxōsus, saxōsa, saxōsum}.) Rough? (\textbf{cōnfragōsus, cōnfragōsa, cōnfragōsum}.) To slide down? (\textbf{dēlābor, dēlābī}.) A river? (\textbf{amnis, amnis}.) A torrent? (\textbf{torrēns, torrentis}.) A confluence? (\textbf{cōnfluēns, cōnfluentis}.) To flow? (\textbf{fluō, fluere}.) To swim across? (\textbf{trānsnatō, trānsnatāre}.) A pool? (\textbf{stagnum, stagnī}.) A whirlpool? (\textbf{gurges, gurgitis}.) A chasm? (\textbf{vorāgō, vorāginis}.) To trickle? (\textbf{mānō, mānāre}.) To flood? (\textbf{inundō, inundāre}.) A beech tree? (\textbf{fāgus, fāgī}.) A fir tree? (\textbf{abiēs, abietis}.) An elm tree? (\textbf{ulmus, ulmī}.) An ash tree? (\textbf{fraxinus, fraxinī}.) Shady? (\textbf{opācus, opāca, opācum}.) A downpour? (\textbf{imber, imbris}.) To be soaking wet? (\textbf{madeō, madēre}.) Rainy? (\textbf{pluviālis, pluviāle}.) Stormy? (\textbf{nimbōsus, nimbōsa, nimbōsum}.) Turbulent? (\textbf{turbidus, turbida, turbidum}.) Hoarfrost? (\textbf{pruīna, pruīnae}.) The cold? (\textbf{frīgus, frīgoris}.) To freeze? (\textbf{gelō, gelāre}.) To be stiff with cold? (\textbf{rigeō, rigēre}.) To feel cold? (\textbf{algeō, algēre}.)
\end{tcolorbox}
